\clearpage
\phantomsection% 
\addcontentsline{toc}{chapter}{RINGKASAN}
\thispagestyle{fancy}

\begin{center}
	\large \bfseries \MakeUppercase{Ringkasan}\\
	\normalsize \normalfont {\thetitle}\\
	\normalsize \normalfont {\theauthor}\\
	\bigskip
	
	\normalsize \normalfont \justifying
	Universitas HKBP Nommensen merupakan salah satu perguruan tinggi swasta yang menyelenggarakan proses Penerimaan Mahasiswa Baru (PMB) setiap tahun ajaran. Proses PMB yang berjalan saat ini masih menggunakan cara manual, di mana calon mahasiswa harus datang langsung ke kampus untuk melakukan pendaftaran, menyerahkan berkas fisik, dan melakukan pembayaran secara konvensional. Kondisi ini menimbulkan berbagai permasalahan, antara lain lambatnya pemrosesan data pendaftar, kesulitan dalam verifikasi berkas, potensi kehilangan data, serta kurangnya transparansi informasi bagi calon mahasiswa mengenai status pendaftaran mereka. Selain itu, pihak panitia PMB mengalami kesulitan dalam mengelola data pendaftar yang berjumlah besar secara efisien.\par

	Berdasarkan permasalahan tersebut, penelitian ini merumuskan dua masalah utama, yaitu: (1) Bagaimana mengembangkan sistem terintegrasi PMB berbasis web yang mampu mengelola seluruh proses penerimaan secara digital? dan (2) Bagaimana menguji fungsionalitas serta kebergunaan sistem yang dikembangkan? Tujuan penelitian ini adalah merancang dan mengembangkan sistem PMB berbasis web yang terintegrasi, serta melakukan pengujian terhadap fungsionalitas dan \textit{usability} sistem tersebut.\par

	Metode pengembangan yang digunakan adalah metode \textit{Prototype}. Metode ini dipilih berdasarkan analisis komparatif terhadap beberapa model pengembangan perangkat lunak seperti \textit{Waterfall}, \textit{Agile}, dan \textit{Spiral}. Metode \textit{Prototype} dinilai paling sesuai karena memungkinkan pengembangan bertahap dengan umpan balik langsung dari pengguna, sehingga sistem yang dihasilkan lebih sesuai dengan kebutuhan aktual. Tahapan pengembangan meliputi identifikasi masalah melalui observasi dan wawancara, pengumpulan data melalui studi pustaka dan kuesioner, analisis kebutuhan, perancangan sistem, implementasi, pengujian, serta evaluasi dan dokumentasi.\par

	Sistem dirancang menggunakan arsitektur \textit{Model-View-Controller} (MVC) dengan \textit{framework} Spring Boot sebagai \textit{backend}, basis data MySQL untuk penyimpanan data, serta antarmuka \textit{responsive} yang dibangun menggunakan HTML, CSS, dan Bootstrap. Fitur-fitur utama sistem meliputi: pendaftaran calon mahasiswa secara daring, unggah dan verifikasi berkas, integrasi pembayaran melalui BRIVA (\textit{Bank Rakyat Indonesia Virtual Account}), notifikasi melalui email menggunakan protokol SMTP, pengelolaan jadwal seleksi, pengumuman hasil seleksi, serta proses daftar ulang secara digital. Sistem juga menyediakan \textit{dashboard} bagi administrator untuk memantau dan mengelola seluruh proses PMB.\par

	Perancangan sistem mencakup pembuatan \textit{use case diagram} untuk empat aktor utama (Admin, Calon Mahasiswa, Pimpinan, dan Panitia PMB), \textit{activity diagram}, \textit{sequence diagram}, \textit{class diagram}, \textit{Data Flow Diagram} (DFD) hingga level detail, \textit{Entity Relationship Diagram} (ERD) dengan sembilan \textit{data store}, serta 15 desain antarmuka pengguna menggunakan Figma. Analisis kebutuhan mengidentifikasi 48 kebutuhan fungsional dan 10 kebutuhan non-fungsional yang harus dipenuhi oleh sistem.\par

	Pengujian sistem direncanakan menggunakan tiga pendekatan, yaitu: (1) \textit{Black-box Testing} dengan 10 \textit{test case} untuk memvalidasi fungsionalitas sistem, (2) \textit{Integration Testing} untuk menguji keterpaduan antar modul dan komponen sistem, serta (3) \textit{System Usability Scale} (SUS) dengan 10 pernyataan skala Likert untuk mengukur tingkat kebergunaan sistem dari perspektif pengguna akhir. Hasil penelitian ini diharapkan menghasilkan sistem PMB yang efisien, terintegrasi, dan mudah digunakan oleh seluruh pemangku kepentingan di Universitas HKBP Nommensen.\par

	\vfill
	
\end{center}
\clearpage
